\section{Relation to the first version}\label{app:changes}

The first published version of this programme~\cite{Sikula2026v1} stated it in a broad form. This
appendix relates the two: what is kept, what is reformulated, and what is no longer claimed. Every change
results from a test or a derivation in the present work.

\subsection{Kept}

\begin{itemize}
\item The primacy of causality over geometry; events without intrinsic properties; the causal order as
  the deepest structure.
\item Causal continuity: influence is conserved at every event.
\item Finite capacity, now as a bound on the flow through an event.
\item The local response variables $\Pi_v$, $I_v$ and $Z_v=I_v/\Pi_v$, and the Causal Resistance
  Functional as the action of a flow.
\item The picture of $c$ as the maximal congestion-free throughput and of $c_{\rm eff}$ falling where the
  substrate is loaded. It is now the throughput principle, which gives $\gamma=1$ in three spatial
  dimensions.
\item The separation of the causal, bridge and observational layers.
\item The ontological themes: observation and causal resolution, branching and merging, horizons as
  limits of reconstructability, superposition as coexisting continuations.
\end{itemize}

\subsection{Reformulated}

\begin{itemize}
\item \textbf{The substrate.} A graph with bounded degree $K_0$, locality by graph distance and time as a
  count of updates has a rest frame. It is replaced by a Lorentz-invariant causal set, as required by the
  theorem of Bombelli, Henson and Sorkin and supported by the tests F0.3, F1.8, F6.1 and F6.2.
\item \textbf{Dimension.} The effective dimension is no longer defined by the growth of graph balls; it is
  read from the order (fraction of related pairs, blind instruments).
\item \textbf{The dynamical principle.} The Causal Resistance Functional is the action of a flow; the
  dynamics of the order is the Benincasa--Dowker action.
\item \textbf{Gravity.} The inverse-square law from flux through spherical shells is replaced by the
  response of the order: clocks, cones and density, with the tensor coupling derived and tested.
\item \textbf{Special relativity.} The Lorentz factor is no longer derived from clocks at rest relative to
  the substrate; Lorentz invariance is a property of the substrate.
\item \textbf{Gravity changes the causal order.} Near matter the cones narrow and events are added;
  causality itself (no loops, conservation of flow) is untouched.
\end{itemize}

\subsection{No longer claimed}

\begin{itemize}
\item The cosmological anchor $a_0=cH_0/2\pi$ and the low-acceleration regime built on it, including the
  reinterpretation of dark matter. The anchor is not derived from the present formulation, and its
  mechanism relied on the bounded-degree substrate.
\item The plateau $D_{\rm eff}\approx3$ as an attractor of a bounded-degree dynamics.
\item Threshold effects in strong gravity (suppression of the redshift of white dwarfs, environment-dependent
  delays of gravitational waves), and strong-field properties of a scalar field equation. Beyond the
  linear order the theory is completed by general relativity.
\item The projection of time and length through the Hubble scale and the projection of impedance through the
  vacuum electromagnetic impedance.
\end{itemize}
